\section{Extended Related Work and Substrate Positioning}
\label{app:extended-related}

This appendix expands Section~\ref{sec:related}.  It is a substrate-positioning
map, not a second contribution section and not an ownership-proof claim.

\paragraph{Generated-text watermarking.}
Token-level watermarking methods mark generated text so that a detector can
identify samples produced under a watermarking rule
\cite{kirchenbauer2023watermark,zhao2023provable,hu2024unbiased,kuditipudi2024robust,christ2023undetectable,lee2024who,chang2024postmark,huang2025waterpool,dabiriaghdam2025simmark,wouters2024optimizing,liu2024adaptive,zhu2024duwak,lau2024waterfall,jiang2025stealthink}.
SynthID-Text is an important production-scale example because it reports a
large deployment and detector evaluation over generated text
\cite{dathathri2024scalable}.  These methods use final text as the evidence
object, making them highly relevant to public and portable detection.  The VSG
result does not claim watermark failure in general.  It shows that the current
public final-text predicates in our artifacts recover zero codeword blocks and
are vulnerable to source-mismatch spoofing.

\paragraph{Publicly auditable watermark protocols.}
Publicly-Detectable Watermarking, Puppy, PVMark, and VOW strengthen the
verification interface by adding public detectability, public verifiability,
commitments, oblivious detection, or zero-knowledge style protocol structure
\cite{fairoze2023publicly,isler2023puppy,duan2025pvmark,luan2026vow}.  These
works are central to the VSG comparison because they show that public audit can
be pursued by changing the evidence substrate.  Our manuscript should not be
read as a claim against those protocol-backed routes.  It asks what remains
when the verifier sees only ordinary final text and no protocol or proof object.

\paragraph{Credential and proof systems.}
C2PA Content Credentials bind assertions, claims, and signatures in a manifest
associated with a digital asset \cite{c2pa2025spec}.  zkLLM and ZKPROV explore
proof-based verification for LLM computation or dataset provenance
\cite{sun2024zkllm,namazi2025zkprov}.  These systems are strong examples of
substrate displacement: they can support deterministic checks when the
credential, statement, proof, or transcript is available.  We do not claim
cryptographic provenance in this manuscript; we use these systems to mark the
difference between proof-bearing substrates and ordinary copied natural text.

\paragraph{Removal, rewriting, and learnability.}
Watermark removal and residual-signal work shows that text-level provenance can
be changed, removed, inherited, or transformed under downstream processing
\cite{chen2025demark,sander2024radioactive}.  Impossibility, learnability, and
spoofing studies further caution that watermark signals can be unstable or
learnable under some attacker models
\cite{zhang2023watermarkssand,gu2023learnability,an2025ditto}.  This
literature motivates our separation between shallow distributional
separability and deterministic codeword recovery.  The current manuscript does
not claim sanitizer robustness or semantic-paraphrase robustness; those would
require separate gates.

\paragraph{Active fingerprints.}
LLM fingerprinting methods use prompts, response patterns, hidden triggers,
human-readable signatures, chained responses, cross-turn correlations, model
identification behavior, probabilistic traces, and scalable verification
\cite{xu2024instructional,zeng2024huref,russinovich2024chainhash,pasquini2025llmmap,yang2024fingerprint,iourovitski2024hide,yamabe2025mergeprint,xu2025ctcc,xu2025evertracer,bai2025esf,hu2025resf,nasery2025scalable,tsai2025rofl}.
These systems are closest to model-side verification.  Our paper should not be
read as a superiority comparison.  It separates the substrate question: does
the evidence live in a cooperative query/trace regime, or does it survive as a
deterministic witness in public final text?

\paragraph{Model-state ownership evidence.}
Classical DNN watermarking and ownership methods embed signatures in
parameters, activations, or trigger responses
\cite{uchida2017embedding,zhang2018protecting,adi2018turning,rouhani2019deepsigns,jia2021entangled,szyller2021dawn,krauss2024clearstamp}.
These substrates can support strong claims under white-box or cooperative
access.  They do not automatically solve the public final-text-only setting in
which a third party has only copied natural text.

\paragraph{Post-release modification.}
Fingerprint erasure, backdoor removal, sleeper-agent persistence, LoRA,
weight averaging, task arithmetic, and model merging all show that downstream
model behavior can change after release
\cite{zhang2025meraser,zeng2024beear,li2024backdoor,hubinger2024sleeper,hu2021lora,wortsman2022model,ilharco2023task,yadav2023ties}.
This motivates scenario-specific evaluation rather than broad robustness
language.  The current VSG manuscript does not claim robustness to arbitrary
fine-tuning, merging, semantic rewriting, or probe-free distillation.

\paragraph{Passive provenance and exposure.}
Passive provenance methods infer exposure, memorization, or lineage without an
owner-controlled active evidence channel
\cite{carlini2024stealing,liang2024model,xie2024recall,puerto2025scaling,ravichander2025imprint}.
They are complementary to the VSG framing.  They can support forensic analysis
in some regimes, but they do not by themselves establish that ordinary final
text carries a deterministic codeword witness.

\begin{table}[t]
\centering
\scriptsize
\setlength{\tabcolsep}{2pt}
\renewcommand{\arraystretch}{1.08}
\begin{tabular}{p{0.18\linewidth}p{0.18\linewidth}p{0.19\linewidth}p{0.18\linewidth}p{0.19\linewidth}}
\toprule
Method family & Evidence substrate & What it can support & Needed assumption & VSG positioning \\
\midrule
Private-key text watermarks & Final generated text plus detector key & Distributional detection or keyed detection & Text remains sufficiently unmodified; detector has key & Closest final-text substrate; current artifacts do not claim watermark success \\
Publicly auditable watermarks & Final text plus public protocol, commitment, or proof object & Publicly auditable detection under protocol assumptions & Protocol artifacts or public verification material are available & Stronger substrate than ordinary final text alone \\
C2PA-style credentials & Asset metadata manifest plus signatures & Manifest integrity and asserted provenance metadata & Credential survives copying and trust anchors are accepted & Metadata substrate; not the copied-text-only regime \\
ZK/proof systems & Statement, witness, and proof transcript & Deterministic proof of specified computation or provenance statement & Statement is well-formed and proof substrate is available & Proof substrate; not claimed by current trace-bound artifacts \\
Provider traces & Token-event logs or first-divergence traces & Cooperative provider-side controllability diagnostics & Authorized verifier receives trace bundle & Supported only as trace-bound diagnostic in current results \\
Model-state fingerprints & Weights, activations, or trigger behavior & White-box or cooperative ownership evidence & Model access or cooperative query conditions & Complementary model-state substrate \\
Public final text only & Ordinary copied natural text & Portable observation; in current artifacts zero codeword recovery & No logs, proofs, metadata, or model access & Current tested predicates are spoofable and not deterministic verification \\
\bottomrule
\end{tabular}
\caption{Related work positioned by evidence substrate.  The table is a claim
scope map, not a ranking and not an ownership-proof claim.  The VSG result is
that trace-bound evidence can be recoverable while the tested public final-text
predicates recover no codeword blocks and can be spoofed by source-mismatch
examples.}
\label{tab:extended-substrate-positioning}
\end{table}

\paragraph{Substrate takeaway.}
Across these families, strong evidence usually depends on the substrate being
available.  The present manuscript's contribution is to make that dependence
explicit: trace-bound first-divergence diagnostics show provider-side event
capacity, while current public final-text predicates do not recover codewords
and can be optimized by source-mismatch examples.
